% === Section 3: Related Work ===
\section{Related Work}\label{sec:related}

The ORIUS framework sits at the intersection of several research threads: runtime
verification, conformal prediction, safety shielding, and CPS telemetry resilience.
We survey each in turn, highlighting the gap that ORIUS addresses.

\subsection{Runtime Verification and Monitoring}

Runtime-verification (RV) monitors check whether an execution trace satisfies a
temporal-logic specification~\cite{leucker2009brief,bartocci2018lectures}.  Classic RV
monitors are \emph{passive}: they flag violations but do not modify the controller's
action.  The simplex architecture~\cite{sha2001using} introduces a \emph{decision
module} that switches between a complex controller and a verified-safe baseline, but
it assumes that the state is accurately known---the observation pipeline is outside
its trust model.

ORIUS goes beyond passive monitoring by actively computing a
constraint-tightened action set and repairing unsafe actions.  Crucially, it operates
on \emph{degraded} observations and inflates the safety envelope accordingly, a
capability absent from standard simplex or watchdog designs.

\subsection{Conformal Prediction for Safety}

Conformal prediction (CP) constructs distribution-free prediction intervals with
finite-sample coverage guarantees~\cite{vovk2005algorithmic}.  Adaptive Conformal
Inference~\cite{gibbs2021adaptive} and related methods handle distribution shift
online.  Recent work applies CP to robotics~\cite{lindemann2023safe} and
autonomous driving~\cite{luo2024conformal_av}.

ORIUS uses CP as the \emph{calibration engine} (Stage~2 of \dcs{}) but wraps it in a
full runtime safety layer with constraint tightening, action shielding, and governance
certification.  Stand-alone CP only produces intervals; ORIUS consumes them,
inflates them according to degradation severity, and feeds the result into a formal
contract stack.

\subsection{Safety Shielding and Control Barrier Functions}

Shield synthesis computes a reactive strategy that enforces a safety specification by
overriding unsafe controller actions~\cite{bloem2015shield}.  Control Barrier
Functions (CBFs)~\cite{ames2019control} certify forward invariance of safe sets.
Both approaches typically require a known, accurate state vector.

ORIUS contributes by coupling the shield to an \emph{observation-consistent state set}
$\Ut$ rather than a point state estimate.  The constraint-tightening stage
(\cref{sec:architecture}) is conceptually analogous to robust CBF formulations, but it
is implemented via interval arithmetic on conformal sets, making it compatible with
black-box learned models.

\subsection{Telemetry Resilience in CPS}

Fault-tolerant control and networked control systems address packet loss and
delay~\cite{hespanha2007survey}.  However, these methods typically assume known
fault models (\eg~Bernoulli packet loss with known rate) and linear dynamics.
ORIUS's detect stage uses statistical detectors (Page--Hinkley, staleness counters)
to \emph{assess} the fault type without assuming a parametric model, and the
calibration stage inflates uncertainty in proportion to the detected severity.

\subsection{Comparison Summary}

\Cref{tab:related_cmp} summarises the coverage of requirements R1--R5 across related
approaches.

\begin{table}[t]
\caption{Comparison of ORIUS with related approaches on
Requirements R1--R5.  \checkmark = addressed, ($\sim$) = partial, $\times$ = not addressed.}
\label{tab:related_cmp}
\centering\small
\begin{tabular}{lccccc}
\toprule
\textbf{Approach} & \textbf{R1} & \textbf{R2} & \textbf{R3} & \textbf{R4} & \textbf{R5} \\
\midrule
Classic RV monitor          & $\sim$ & $\times$ & $\times$    & $\times$ & $\times$ \\
Simplex architecture        & $\times$ & $\times$ & $\times$  & \checkmark & $\times$ \\
Conformal prediction        & $\times$ & \checkmark & $\times$ & $\times$ & $\times$ \\
CBF / Shield synthesis      & $\times$ & $\times$ & \checkmark & \checkmark & $\times$ \\
Fault-tolerant NCS          & \checkmark & $\sim$ & $\times$   & $\times$ & $\times$ \\
\midrule
\textbf{ORIUS (\dcs{})}    & \checkmark & \checkmark & \checkmark & \checkmark & \checkmark \\
\bottomrule
\end{tabular}
\end{table}

\subsection{Positioning}

ORIUS is, to our knowledge, the first runtime layer that satisfies all five
requirements simultaneously while remaining domain-agnostic through thin typed
adapters.  The two-domain validation we present (\cref{sec:battery,sec:av})
demonstrates that the same kernel---compiled once, configured via adapter---achieves
zero true-state violations on battery dispatch and closes the OASG on autonomous
vehicle prediction.
